% Page 029 — page imprimée 25
% Contenu seul : le préambule et la pagination sont gérés par meyrueis.tex

\begin{center}
{\Large\bfseries FABRIQUE DE POINTES\\
A MEYRUEIS}

1835\#
\end{center}

M. Felix Laporte sous la raison sociale Laporte et Compagnie, a établi à Meyrueis,
(\textit{non loin de Lucallous semble-t-il}) une fabrique de pointes, dites de Paris,
de becquets ou clous à vis et d’aiguilles à tricoter.

Cette fabrique, introduite depuis quelques années, à Meyrueis par plusieurs actionnaires, n’a
pris une direction et une activité convenables qu’à compter de la fin de 1832, époque à
laquelle M. Laporte a succédé aux anciens propriétaires.

D’après un exposé certifié par le maire de cette ville, elle occupe de 300 à 350 ouvriers,
hommes, femme et enfants, et assure ainsi du pain à une partie de la population, auparavant
inactive.

Elle tire la matière première du département du Doubs et si les forges d’Alais peuvent tréfiler
le fer à un point convenable, elle sera à même de payer les frais de transport à 80\% meilleur
marché et de diminuer d’autant le prix de vente. Ses débouchés sont dans les départements
méridionaux.

Quant à ses prix de vente ils sont modérés et la fabrication des produits est de bonne
qualité ... Dans un pays pauvre comme le nôtre on doit des encouragements à toutes les
importations industrielles...\footnote{Revue de la Société d’Agriculture de la ville de Mende Lozère 1833}

\begin{quote}
\itshape
En 1834, lors d’une exposition des produits de l’Industrie à Mende ? Paris ? Il a été accordé une citation à M.
Laporte et Cie pour la fabrication à Meyrueis de « pointes, clous et aiguilles à bas ».
\end{quote}

\begin{flushright}
La France Pittoresque (Lozère) 1835
\end{flushright}

\vspace{1.2em}

\begin{center}
{\large\bfseries Mémoire sur la ville de Meyrueis \quad (1881)}
\end{center}

« ...C’est une contrée intéressante : la cinquième ville du Diocèse Elle est habitée par une
noblesse et une bourgeoisie nombreuse. On y compte en ce moment douze Chevaliers de St.
Louis et beaucoup d’autres militaires au service du Roi ; beaucoup de bourgeois vivant
noblement, plusieurs avocats ou gradués ; quatre offices de notaire, ses fabriques d’étoffe de
laine de son cru qui partent presque toutes pour l’étranger ne pouvant être servies que par ses
habitants.

Le commerce qui s’y fait des plus beaux mulets du Royaume et des meilleurs moutons la rend
pour ainsi dire, l’entrepôt de la Provence, des provinces voisines et même de l’Espagne.
C’est la mère nourrice des Cévennes par le blé qu’elle leur fournit.
...Il faudrait supprimer toutes les justices bannerettes... que les places de judicature ne
fussent plus vénales, que ce fut le mérite seul qui put y appeler et que le tarif le plus exact
réglât tous les droits ».

\begin{quote}
\itshape
Mémoire sur la ville de Meyrueis, (Annuaire de la Lozère 1881) ainsi que la prise du Malzieu
le 17 Novembre 1573 Pages 3 et ss. \quad Notices rédigées par Fernand André, Archiviste.
\end{quote}
